\section{Datasets, Splits, and Preprocessing}
\label{app:datasets}

\paragraph{FACED.} 123 subjects, 28 emotional video clips covering 9
emotion categories (Anger, Disgust, Fear, Sadness, Neutral,
Amusement, Inspiration, Joy, Tenderness) at 3--4 stimuli per class.
EEG: 32 channels, 250\,Hz sampling, 30-second clip duration. We use
the official preprocessing protocol distributed with the dataset
release \citep{chen2023faced}: bandpass 0.5--47\,Hz, notch 50\,Hz,
ICA-based artefact rejection, common-average re-reference. Train /
validation / test split: subjects 0--79 / 80--99 / 100--122 (no
subject overlap), matching the protocol used by CBraMod
\citep{wang2025cbramod} and EMOD \citep{chen2025emod} for fair
comparison.

\paragraph{SEED-V.} 5-class emotion (Disgust, Fear, Sadness, Neutral,
Happy), used for cross-dataset transfer validation only. Standard
SEED-V split with subject-disjoint train / test, 62-channel EEG at
1000\,Hz downsampled to 200\,Hz. We replicate the SEED-V CBraMod
recipe at $d{=}6$ for the ensemble-generality numbers reported in
Section~\ref{sec:sota}.

\paragraph{DE features.} For each (subject, stimulus, channel) we
compute differential entropy in five canonical bands
($\delta$ 0.5--4, $\theta$ 4--8, $\alpha$ 8--13, $\beta$ 13--30,
$\gamma$ 30--47\,Hz) per second over the 30-second clip, then mean
over time to obtain a $32 \times 5$ feature matrix per
(subject, stim) pair. Cohort-level analyses average across the 123
subjects to give the $28 \times 32 \times 5$ tensor used in
Section~\ref{sec:vaxis-brain}.

\section{V-Axis Extraction Protocol}
\label{app:vaxis-protocol}

\paragraph{Story prompts.} Nine emotion-evocative stories (one per
FACED class), each 1--3 sentences. Stories were authored once,
reviewed by two annotators blind to the LLM evaluation results, and
held fixed across all 18 LLMs. Full prompt list in supplementary
materials.

\paragraph{Paraphrase generation.} For each story, we generate $n=50$
paraphrases via independent calls to a generic paraphrasing LLM
(Qwen3.5-7B with temperature 0.7), with the prompt: "Rewrite the
following so it preserves all emotional content but uses different
words and sentence structures: \{story\}". Paraphrases are filtered
for length (5--80 tokens) and minimum cosine distance from the
source ($\geq 0.10$ on Sentence-BERT embeddings).

\paragraph{Hidden-state extraction.} Each paraphrase is tokenised and
passed through the target LM. We take the last-token hidden state at
layer $L$ (default: $L_{\max}-1$, i.e.\ the penultimate layer; for
Qwen3.5-1.7B this is $L=27$ of 28). Class centroids are
$c_k = \frac{1}{50} \sum_{i=1}^{50} h_{k,i}^{(L)}$ for $k=1,\dots,9$.

\paragraph{V-axis as PCA-PC1.} Stack the nine centroids into a
$9 \times D$ matrix and run mean-centred PCA. The V-axis is the unit
vector along PC1, oriented so that $\langle c_{\mathrm{Joy}},
v\rangle > 0$. PC1 explains 41--58\% of the across-class variance
across the 18 models we tested (median 49\%).

\paragraph{Robustness.} Robustness to story prompt rephrasing was
verified by sampling 5 alternative prompt sets per class and
comparing per-stimulus V-axis projections; mean cross-prompt
correlation $r > 0.93$ (across 18 models). Robustness to paraphrase
count was verified at $n \in \{1, 5, 25, 50\}$; SST-2 zero-shot AUC
of $0.74$ at $n=1$, saturating to $0.868$ at $n=5$.

\section{EEG Model Training Details}
\label{app:training}

\paragraph{Architecture.} EMOD axial transformer
\citep{chen2025emod}: input $32 \times 5$ DE features, axial
self-attention over channels and bands at depth $d \in \{3, 6\}$,
filter dimension $f=128$. The full SOTA recipe is $d=6$, $f=128$.

\paragraph{Optimisation.} AdamW with $\mathrm{lr}=10^{-3}$,
$\beta_1=0.9$, $\beta_2=0.999$, weight decay $10^{-2}$. Cosine
schedule with 5-epoch linear warmup. Batch size $128$. Training
length $e \in \{100, 150\}$ epochs. We use the validation BACC for
checkpoint selection and report test BACC.

\paragraph{Augmentation.} Per-trial Gaussian noise
($\sigma=0.05 \cdot \mathrm{std}$ of feature distribution), random
channel dropout ($p=0.15$), and random temporal masking
($\mathrm{up\ to\ }5\%$ of seconds). All augmentations are applied
with $p=0.6$ during training only.

\paragraph{Knowledge distillation.} Soft-target KD with rand9 9-D
orthonormal teacher (no LLM content); $\lambda_{\mathrm{KD}} = 0.5$,
$T = 1.0$. The LLM-9-D and rand9-9-D variants give within-noise
identical BACC, validating that KD provides architectural
regularisation without semantic content
(Section~\ref{sec:sota-component-ablation}).

\paragraph{Ensemble.} 10 checkpoints split as 5 seeds at $e=100$ +
5 seeds at $e=150$. Seeds: $\{42, 123, 456, 789, 2025\}$ each.
Ensemble prediction is uniform softmax averaging across the 10
checkpoints' output probabilities, followed by argmax.

\section{Statistical Methods}
\label{app:statistics}

\paragraph{Bootstrap CIs.} Per-subject and per-stim 95\% CIs are
computed from $B=10{,}000$ subject-resampled bootstraps. For cohort
correlations, we report the bootstrap median, 2.5\% / 97.5\%
percentile interval, and a Fisher-$z$ $p$-value for $r=0$.

\paragraph{Random-direction null.} For cohort EEG--LLM correlations,
we sample $N=200$ random Gaussian directions in $\mathbb{R}^{28}$
matched in $L_2$ norm to the V-axis projection, recompute the cohort
$|r|$ for each, and report the percentile of the empirical $|r|$
in the null distribution.

\paragraph{Cross-architecture null.} For the 36-checkpoint cross-arch
correlation, we sample $N=100$ random directions matched in norm to
$v_{\mathrm{CLIP}}$, and recompute the BACC--V-axis correlation. The
null is distributed broadly ($\sigma=0.62$) due to the low intrinsic
dimension of the class-PC1 subspace; we report the empirical
correlation's percentile rank.

\paragraph{Per-checkpoint LOO contribution.} For the 10-vanilla
$d{=}6$ ensemble, we compute
$\Delta_k = \mathrm{BACC}(\mathrm{ens}) -
\mathrm{BACC}(\mathrm{ens}\setminus k)$ for each $k$. Reported
correlations between $\Delta_k$ and per-ckpt features are
Pearson with $n=10$.

\paragraph{V-axis intervention comparisons.} For each intervention,
$\Delta$ is computed seed-paired with the matched-recipe vanilla
baseline; $p$-values are paired $t$-tests across 5 seeds.

\section{Full V-Axis Intervention Table}
\label{app:full-vaxis-table}

\begin{table}[h]
\centering
\small
\begin{tabular}{llrrl}
\toprule
\# & Family / variant & $\Delta$ & $p$ & Verdict \\
\midrule
1  & Frontal-mask $\lambda{=}0.5$        & $-0.052$ & $0.0015$ & sig.\ negative \\
2  & Frontal-mask $\lambda{=}0.1$        & $-0.009$ & $0.21$  & n.s. \\
3  & FAA $\lambda{=}0.5$                 & $-0.044$ & $0.006$ & sig.\ negative \\
4  & FAA $\lambda{=}0.1$                 & $-0.018$ & $0.063$ & borderline \\
5  & Anger-weighted $\lambda{=}0.5$      & $-0.054$ & $0.0003$ & sig.\ negative \\
6  & Occipital $\lambda{=}0.1$           & $-0.022$ & $0.007$ & sig.\ negative \\
7  & Topo-optimal $\lambda{=}0.1$        & $-0.013$ & $0.039$ & sig.\ negative \\
8  & Topo-optimal $\lambda{=}0.05$       & $+0.0021$ & $0.50$ & n.s. (NULL) \\
9  & Procrustes $\lambda{=}0.05$         & $+0.0012$ & $0.89$ & n.s. (NULL) \\
10 & RSA $\lambda{=}5.0$                 & $-0.093$ & $<10^{-4}$ & sig.\ negative \\
11 & RSA $\lambda{=}1.0$                 & $-0.057$ & $<10^{-3}$ & sig.\ negative \\
12 & Distance-CE $\tau{=}5.0$            & $-0.397$ & $<10^{-4}$ & catastrophic \\
13 & Multi-V $\lambda{=}0.5$             & $-0.043$ & $<10^{-3}$ & sig.\ negative \\
14 & EEG-AUX MSE $\lambda{=}0.5$         & $-0.043$ & $<10^{-3}$ & sig.\ negative \\
15 & EEG-AUX MSE $\lambda{=}0.1$         & $-0.016$ & $0.04$  & sig.\ negative \\
16 & EMODSTYLE class $\lambda{=}1.0$     & $-0.010$ & $0.18$  & n.s. \\
17 & EMODSTYLE stim $\lambda{=}0.5$      & $+0.0066$ & $0.22$ & n.s. (sweet spot) \\
18 & PEFT fullhead $\lambda{=}0.1$       & $-0.009$ & $0.069$ & borderline \\
19 & PEFT LoRA $\lambda{=}0.1$           & $-0.010$ & $0.12$  & n.s. \\
20 & PEFT IA3 $\lambda{=}0.1$            & $-0.016$ & $0.05$  & sig.\ negative \\
21 & Pretrain-FT (frozen)                & $-0.401$ & $<10^{-4}$ & catastrophic \\
22 & Pretrain-FT (unfrozen)              & $-0.190$ & $<10^{-4}$ & sig.\ negative \\
23 & XEEG (FACED $\to$ SEED-V)           & $+0.0015$ & $0.97$  & n.s. (zero) \\
24 & SCALING (full data) $\lambda{=}0.5$ & $-0.045$ & $<0.01$  & sig.\ negative \\
25 & Uncertainty multi-task              & $-0.067$ & $<10^{-3}$ & sig.\ negative \\
\midrule
   & SOTA recipe + Topo $\lambda{=}0.1$  & $-0.015$ & $0.020$  & sig.\ negative (cliff) \\
   & SOTA recipe + EMODSTYLE             & $-0.024$ & $0.001$  & sig.\ negative (cliff) \\
\bottomrule
\end{tabular}
\caption{Full V-axis-as-supervision intervention results. 13 / 25
families produce a statistically significant negative; the remaining
12 are not significantly different from zero. None reach $p<0.05$ in
the positive direction at the strong baseline. Below the line are
the SOTA-recipe replications confirming the saturation cliff.}
\label{tab:full-vaxis}
\end{table}

\begin{figure}[ht]
\centering
\IfFileExists{figures/landmark/lf9_all_interventions_forest.pdf}
  {\includegraphics[width=0.9\linewidth]{landmark/lf9_all_interventions_forest.pdf}}
  {\includegraphics[width=0.9\linewidth]{F9_all_interventions.pdf}}
\caption{Forest plot of all V-axis-as-supervision interventions
(Table~\ref{tab:full-vaxis}) ranked by $\Delta$ vs.\ baseline, with
significance tier coloured. None of the 25 reach $p<0.05$ in the
positive direction at the strong baseline.}
\label{fig:forest}
\end{figure}

\section{FACED Test Confusion Matrices}
\label{app:confusion}

\begin{figure}[ht]
\centering
\IfFileExists{figures/landmark/lf10_confusion_matrices.pdf}
  {\includegraphics[width=0.95\linewidth]{landmark/lf10_confusion_matrices.pdf}}
  {\includegraphics[width=0.95\linewidth]{F10_two_tier_schematic.pdf}}
\caption{FACED 9-class test confusion. Left: best single checkpoint
(BACC $0.6755$). Right: 10-checkpoint ensemble (BACC $0.6948$). The
ensemble's gain is concentrated on the off-diagonal cells where
single-seed disagreement is highest, consistent with the within-class
residual mechanism of Section~\ref{sec:two-tier}.}
\label{fig:confusion}
\end{figure}

\section{Concept-Library Transfer Matrix}
\label{app:concept-library-full}

The full $20\!\times\!20$ cross-concept transfer matrix (each row
is the V-axis of source concept $i$ applied to the pole-vs-pole
benchmark of target concept $j$, best-orientation AUC reported) is
shown in Figure~\ref{fig:concept-library-full}. Rows are sorted by
mean off-diagonal AUC (most-general source axes at the top); the
diagonal is the self-AUC for that concept's own benchmark.

\begin{figure}[h]
\centering
\IfFileExists{figures/landmark/lf15_concept_transfer_heatmap.pdf}
  {\includegraphics[width=\linewidth]{landmark/lf15_concept_transfer_heatmap.pdf}}
  {\includegraphics[width=\linewidth]{landmark/lf15_concept_transfer_heatmap.png}}
\caption{Cross-concept transfer matrix (best-orientation AUC).
Diagonal: self-classification (white-bordered cells). Off-diagonal:
source-axis-applied-to-target-benchmark AUC. Sorted top-to-bottom by
mean off-diagonal AUC, separating \emph{general} concept axes (top:
sarcasm, fear, envy, shame) from \emph{specific} ones (bottom:
complexity, decisiveness). Right marginal: per-source mean
off-diagonal AUC (generality). Top marginal: per-target mean
off-diagonal AUC (captured-by-others). Mean diagonal $0.97$, mean
off-diagonal $0.83$, gap $+0.13$ confirms concept specificity in the
recipe: if every concept axis collapsed onto a single global affect
direction, the off-diagonal would equal the diagonal.}
\label{fig:concept-library-full}
\end{figure}

\section{Negative Results Catalogue}
\label{app:negatives}

In the spirit of full disclosure, the following null or negative
results bear mention beyond what was incorporated into the main
paper:
\begin{enumerate}
\item \textbf{Theta-gamma cross-frequency coupling does not mediate
the V-axis} (Section~\ref{sec:topography}): Tort modulation index
$\rho(\mathrm{PAC}, \mathrm{V\text{-}axis\ }r) = +0.082$, $p=0.667$.
\item \textbf{Mutual-information control is not significant}
(Section~\ref{sec:vaxis-brain}): MI $p=0.115$ vs.\ linear $p=0.020$.
The V-axis-EEG relationship is essentially linear.
\item \textbf{Random-direction null on cross-arch correlation}
(Section~\ref{sec:convergence}): $r=+0.885$ at the
$93^{\mathrm{rd}}$ percentile, $p_{\mathrm{one}}=0.079$. Within-class
residual ($r=+0.74$, $p=0.014$) is the robust signal.
\item \textbf{Per-subject V-axis adaptation does not transfer.} The
per-subject best-channel oracle reaches $|r|=0.62$, but no clean
estimator (V1 global top-K, V2 per-val majority, V3 TTA label-free)
transfers more than $+0.001$ over the cohort top-K.
\item \textbf{Cross-architecture ensembling does not add diversity.}
Mixing $d \in \{4, 6, 8, 10\}$ checkpoints gives within-arch
$0.6798$ vs.\ cross-arch $0.6791$ (n.s.). The diversity is
entirely intra-architecture, training-length-driven.
\item \textbf{Mega-ensembles plateau at 10 checkpoints.} 15-, 20-,
25-checkpoint pools all hit BACC $\leq 0.6948$.
\item \textbf{Stimulus-aggregation re-ranking does not help.}
Operating at the 28-stimulus level instead of trial level gives a
BACC=1.0 by construction (label leak); the corresponding
trial-level head does not transfer.
\item \textbf{Test-time augmentation hurts at K=5} ($0.6753 \to
0.6720$). Averaging predictions over augmented test trials does
not preserve the V-axis residual structure.
\end{enumerate}

\section{Reproducibility}
\label{app:repro}

Code, configs, model checkpoints, and figure-generation scripts will
be released at the camera-ready stage under an MIT licence. Each
experiment in this paper is logged with its SLURM job ID, conda
environment hash, git SHA, and seed list. The 10-checkpoint ensemble
SOTA result is reproducible from a single SLURM array job
(\verb|slurm_d6_e100_e150_5seeds.sh|, $\sim$10 GPU-hours per seed
on V100).

\section{Resource Estimate}
\label{app:compute}

The full paper used approximately $4{,}500$ GPU-hours on V100 nodes:
$\sim 600$ for the recipe ablation cascade, $\sim 1{,}200$ for the
25 V-axis-supervision interventions, $\sim 800$ for the 36-checkpoint
cross-architecture analysis, $\sim 1{,}200$ for ensemble construction
and val-test rank diagnostics, and $\sim 700$ for cross-dataset
generality experiments. The V-axis extraction itself is CPU-cheap
($\sim 10$ minutes per LM on a single 80GB GPU).
